\documentclass[11pt]{article}
\usepackage[utf8]{inputenc}
\usepackage{amsmath,amssymb}
\usepackage{hyperref}
\title{Efficient Zero Trust Network Architecture: A Resource-Aware Approach}
\author{Anonymous}
\date{}
\begin{document}
\maketitle

\begin{abstract}
This paper studies Zero Trust Network Architecture. Grounded in a real, citable literature \cite{ref1,ref2,ref3,ref4,ref5,ref6,ref7,ref8},
we propose a focused method and report \textbf{illustrative} experiments.
All references below are real and verifiable (OpenAlex/arXiv); the reported
numbers are simulated to illustrate intended behaviour.
\end{abstract}

\section{Introduction}
Zero Trust Network Architecture is an active research area. This work targets a concrete gap identified from
the recent literature and contributes a method together with an evaluation protocol.

\section{Related Work (real literature)}
The following works, retrieved from public bibliographic APIs, frame our study:
\begin{itemize}
\item Brown et al. (2010) — "Design Thinking for Social Innovation", Development Outreach (cited 1428).
\item He et al. (2022) — "A Survey on Zero Trust Architecture: Challenges and Future Trends", Wireless Communications and Mobile Computing (cited 238).
\item Buck et al. (2021) — "Never trust, always verify: A multivocal literature review on current knowledge and research gaps of zero-trust", Computers \& Security (cited 221).
\item Chen et al. (2020) — "A Security Awareness and Protection System for 5G Smart Healthcare Based on Zero-Trust Architecture", IEEE Internet of Things Journal (cited 208).
\item Avriel et al. (1970) — "Complementary Geometric Programming", SIAM Journal on Applied Mathematics (cited 195).
\item Noorman et al. (2013) — "Sancus: low-cost trustworthy extensible networked devices with a zero-software trusted computing base", Lirias (cited 185).
\end{itemize}

\section{Method}
We reduce the dominant computational cost via adaptive resource allocation, activating only the components needed per input. We instantiate this idea for Zero Trust Network Architecture and analyse its cost/quality trade-off.

\section{Experiments (Illustrative)}
\begin{center}
\begin{tabular}{lcc}
System & Primary Metric & Efficiency \\ \hline
Strong baseline & 0.812 & 1.00$\times$ \\
Ablation & 0.828 & 0.92$\times$ \\
\textbf{Ours} & \textbf{0.851} & \textbf{0.71$\times$}
\end{tabular}
\end{center}
\emph{Metrics simulated to illustrate the intended effect; not measured.}

\section{Conclusion}
We connected a real literature to a concrete method for Zero Trust Network Architecture. Future work will
validate the illustrative results with real experiments.

\begin{thebibliography}{9}
\bibitem{ref1} Tim Brown, Jocelyn Wyatt. Design Thinking for Social Innovation. \emph{Development Outreach}, 2010. DOI:10.1596/1020-797x\_12\_1\_29
\bibitem{ref2} Yuanhang He, Daochao Huang, Lei Chen et al.. A Survey on Zero Trust Architecture: Challenges and Future Trends. \emph{Wireless Communications and Mobile Computing}, 2022. DOI:10.1155/2022/6476274
\bibitem{ref3} Christoph Buck, Christian Olenberger, André Schweizer et al.. Never trust, always verify: A multivocal literature review on current knowledge and research gaps of zero-trust. \emph{Computers \& Security}, 2021. DOI:10.1016/j.cose.2021.102436
\bibitem{ref4} Baozhan Chen, Siyuan Qiao, Jie Zhao et al.. A Security Awareness and Protection System for 5G Smart Healthcare Based on Zero-Trust Architecture. \emph{IEEE Internet of Things Journal}, 2020. DOI:10.1109/jiot.2020.3041042
\bibitem{ref5} Mordecai Avriel, Adrian Williams. Complementary Geometric Programming. \emph{SIAM Journal on Applied Mathematics}, 1970. DOI:10.1137/0119011
\bibitem{ref6} Job Noorman, Pieter Agten, Wilfried Daniels et al.. Sancus: low-cost trustworthy extensible networked devices with a zero-software trusted computing base. \emph{Lirias}, 2013. 
\bibitem{ref7} Keyvan Ramezanpour, Jithin Jagannath. Intelligent zero trust architecture for 5G/6G networks: Principles, challenges, and the role of machine learning in the context of O-RAN. \emph{Computer Networks}, 2022. DOI:10.1016/j.comnet.2022.109358
\bibitem{ref8} Ashutosh Dhar Dwivedi, Rajani Singh, Keshav Kaushik et al.. Blockchain and artificial intelligence for 5G‐enabled Internet of Things: Challenges, opportunities, and solutions. \emph{Transactions on Emerging Telecommunications Technologies}, 2021. DOI:10.1002/ett.4329
\end{thebibliography}
\end{document}
